\documentclass{article}
\usepackage[utf8]{inputenc}
\usepackage[german]{babel}
\usepackage{amsmath, amssymb, amsthm}
% \usepackage{mathrsfs}
\usepackage[scr]{rsfso}
\usepackage{booktabs}
\usepackage{array}

\newcolumntype{L}[1]{>{\raggedright\arraybackslash}p{#1}}

\title{Beweis von T1 und T3 im Scratch-Format}
\author{Nach Daniel J. Velleman: \emph{How to Prove It}}
\date{\today}

\begin{document}

\maketitle

\section*{1. Definition einer Topologie}

Sei $X$ eine Grundmenge. Ein Mengensystem $\mathcal{T} \subseteq \mathscr{P}(X)$ heißt \textbf{Topologie} auf $X$, wenn die folgenden drei Axiome erfüllt sind:
\begin{enumerate}
    \item T1 (Leere Menge und Ganzraum): $\emptyset \in \mathcal{T} \land X \in \mathcal{T}$
    \item T2 (Endliche Durchschnitte): $\forall U, V \, \big( (U \in \mathcal{T} \land V \in \mathcal{T}) \rightarrow (U \cap V \in \mathcal{T}) \big)$
    \item T3 (Beliebige Vereinigungen): $\forall \mathcal{A} \, \big( \mathcal{A} \subseteq \mathcal{T} \rightarrow \bigcup \mathcal{A} \in \mathcal{T} \big)$
\end{enumerate}

---

\section*{Setup und Voraussetzungen}
Zur Erinnerung: Wir haben eine Familie von Topologien $(\mathcal{T}_i)_{i \in I}$ auf einer Menge $X$. Unser Durchschnittssystem ist definiert als:
\[
\mathcal{T}_{\cap} = \bigcap_{i \in I} \mathcal{T}_i = \{ U \subseteq X \mid \forall i \in I \, (U \in \mathcal{T}_i) \}
\]
Das bedeutet für unsere Beweise: Eine Menge liegt genau dann in $\mathcal{T}_{\cap}$, wenn sie in \emph{jedem einzelnen} $\mathcal{T}_i$ liegt.

---

\newpage

\section*{1. Scratch Work für Axiom T1}

Da im Goal ein UND ($\land$) steht, teilt sich das Rezept in zwei unabhängige Teilziele auf. Wir zeigen hier beispielhaft den Beweis für die leere Menge $\emptyset$.

\begin{center}
\small
\begin{tabular}{ L{0.45\textwidth} | L{0.45\textwidth} }
\toprule
\textbf{Givens} & \textbf{Goals} \\
\midrule
% Schritt 0
1. $\forall i \in I \, (\mathcal{T}_i \text{ ist Topologie})$ \newline
2. $\mathcal{T}_{\cap} = \{ U \subseteq X \mid \forall i \in I \, (U \in \mathcal{T}_i) \}$ &
\emph{Goal:} $\emptyset \in \mathcal{T}_{\cap}$ \newline
\newline
(\emph{Umgewandelt via Definition von $\mathcal{T}_{\cap}$}) \newline
\emph{Goal*:} $\forall i \in I \, (\emptyset \in \mathcal{T}_i)$ \\[1ex]
\hline

% Schritt 1: Allquantor im Goal abbauen
1. [Givens wie oben] \newline
3. Sei $i_0 \in I$ ein beliebiges Indexelement. &
\emph{Neues Goal:} $\emptyset \in \mathcal{T}_{i_0}$ \\[1ex]
\hline

% Schritt 2: Eigenschaft der Familie nutzen
1. [Givens wie oben] \newline
3. $i_0 \in I$ ist beliebig. \newline
4. $\mathcal{T}_{i_0}$ erfüllt Axiom T1 (da es eine Topologie ist!) \newline
5. $\emptyset \in \mathcal{T}_{i_0}$ (Direkt aus Given 4) &
\emph{Goal:} $\emptyset \in \mathcal{T}_{i_0}$ \newline
\newline
\checkmark \quad (Ziel erreicht!) \\
\bottomrule
\end{tabular}
\end{center}

Hier die Schritte fuer $X$

\begin{center}
\small
\begin{tabular}{ L{0.45\textwidth} | L{0.45\textwidth} }
\toprule
\textbf{Givens} & \textbf{Goals} \\
\midrule
% Schritt 0
1. $\forall i \in I \, (\mathcal{T}_i \text{ ist Topologie})$ \newline
2. $\mathcal{T}_{\cap} = \{ U \subseteq X \mid \forall i \in I \, (U \in \mathcal{T}_i) \}$ &
\emph{Goal:} $X \in \mathcal{T}_{\cap}$ \newline
\newline
(\emph{Umgewandelt via Definition von $\mathcal{T}_{\cap}$}) \newline
\emph{Goal*:} $\forall i \in I \, (X \in \mathcal{T}_i)$ \\[1ex]
\hline

% Schritt 1: Allquantor im Goal abbauen
1. [Givens wie oben] \newline
3. Sei $i_0 \in I$ ein beliebiges Indexelement. &
\emph{Neues Goal:} $X \in \mathcal{T}_{i_0}$ \\[1ex]
\hline

% Schritt 2: Eigenschaft der Familie nutzen
1. [Givens wie oben] \newline
3. $i_0 \in I$ ist beliebig. \newline
4. $\mathcal{T}_{i_0}$ erfüllt Axiom T1 (da es eine Topologie ist!) \newline
5. $X \in \mathcal{T}_{i_0}$ (Direkt aus Given 4) &
\emph{Goal:} $X \in \mathcal{T}_{i_0}$ \newline
\newline
\checkmark \quad (Ziel erreicht!) \\
\bottomrule
\end{tabular}
\end{center}

\newpage

\section*{2. Scratch Work für Axiom T2}

Um zu beweisen, dass $\mathcal{T}_{\cap}$ eine Topologie ist, müssen wir alle drei Axiome zeigen. Schauen wir uns das Velleman-Scratch-Work für das oft kniffligste Axiom an: **T2 (Stabilität unter endlichen Durchschnitten)**.

\begin{center}
\small
\begin{tabular}{ L{0.45\textwidth} | L{0.45\textwidth} }
\toprule
\textbf{Givens} & \textbf{Goals} \\
\midrule
% Schritt 0: Ausgangssituation
1. $\forall i \in I \, (\mathcal{T}_i \text{ ist Topologie})$ \newline
2. $\mathcal{T}_{\cap} = \{ U \subseteq X \mid \forall i \in I \, (U \in \mathcal{T}_i) \}$ &
\emph{Goal:} $\forall U, V \big( (U \in \mathcal{T}_{\cap} \land V \in \mathcal{T}_{\cap}) \rightarrow U \cap V \in \mathcal{T}_{\cap} \big)$ \\[1ex]
\hline

% Schritt 1: Goal abbauen (Allquantoren und Implikation)
1. [Givens wie oben] \newline
3. Seien $U$ und $V$ beliebige Mengen. \newline
4. $U \in \mathcal{T}_{\cap}$ (Annahme) \newline
5. $V \in \mathcal{T}_{\cap}$ (Annahme) &
\emph{Neues Goal:} $U \cap V \in \mathcal{T}_{\cap}$ \newline
\newline
(\emph{Laut Definition von $\mathcal{T}_{\cap}$ bedeutet das:}) \newline
\emph{Goal*:} $\forall i \in I \, (U \cap V \in \mathcal{T}_i)$ \\[1ex]
\hline

% Schritt 2: Givens 4 und 5 mittels Definition von T_neu aufspalten
1. [Givens wie oben] \newline
3. $U, V$ beliebig \newline
6. $\forall i \in I \, (U \in \mathcal{T}_i)$ \quad (aus Given 4) \newline
7. $\forall i \in I \, (V \in \mathcal{T}_i)$ \quad (aus Given 5) &
\emph{Goal*:} $\forall i \in I \, (U \cap V \in \mathcal{T}_i)$ \\[1ex]
\hline

% Schritt 3: Allquantor im Goal angreifen
1. [Givens wie oben] \newline
3. $U, V$ beliebig; 6. und 7. wie oben \newline
8. Sei $i_0 \in I$ ein beliebiges, festes Indexelement. &
\emph{Neues Goal:} $U \cap V \in \mathcal{T}_{i_0}$ \\[1ex]
\hline

% Schritt 4: Givens mit dem konkreten i_0 instanziieren
3. $U, V$ beliebig; 8. $i_0 \in I$ beliebig \newline
9. $U \in \mathcal{T}_{i_0}$ \quad (Instanziierung von 6 mit $i_0$) \newline
10. $V \in \mathcal{T}_{i_0}$ \quad (Instanziierung von 7 mit $i_0$) \newline
11. $\mathcal{T}_{i_0}$ erfüllt T2 (da es eine Topologie ist!) &
\emph{Goal:} $U \cap V \in \mathcal{T}_{i_0}$ \\[1ex]
\hline

% Schritt 5: Axiom T2 von T_{i_0} nutzen (Modus Ponens)
% Wir instanziieren das Axiom T2 von T_{i_0} mit unseren Mengen U und V
3. $U,V$ beliebig; 8. $i_0$ beliebig; 9., 10. wie oben \newline
12. $(U \in \mathcal{T}_{i_0} \land V \in \mathcal{T}_{i_0}) \rightarrow U \cap V \in \mathcal{T}_{i_0}$ \newline
13. $U \cap V \in \mathcal{T}_{i_0}$ \quad (Modus Ponens aus 9, 10 und 12) &
\emph{Goal:} $U \cap V \in \mathcal{T}_{i_0}$ \newline
\newline
\checkmark \quad (Ziel erreicht!) \\
\bottomrule
\end{tabular}
\end{center}

\newpage

\section*{3. Scratch Work für Axiom T3 (Beliebige Vereinigungen)}

Das Axiom fordert: Wenn ein Teilsystem $\mathcal{A}$ eine Teilmenge von $\mathcal{T}_{\cap}$ ist, dann muss auch die große Vereinigung $\bigcup \mathcal{A}$ in $\mathcal{T}_{\cap}$ liegen.

\begin{center}
\small
\begin{tabular}{ L{0.45\textwidth} | L{0.45\textwidth} }
\toprule
\textbf{Givens} & \textbf{Goals} \\
\midrule
% Schritt 0
1. $\forall i \in I \, (\mathcal{T}_i \text{ ist Topologie})$ \newline
2. $\mathcal{T}_{\cap} = \{ U \subseteq X \mid \forall i \in I \, (U \in \mathcal{T}_i) \}$ &
\emph{Goal:} $\forall \mathcal{A} \, \big( \mathcal{A} \subseteq \mathcal{T}_{\cap} \rightarrow \bigcup \mathcal{A} \in \mathcal{T}_{\cap} \big)$ \\[1ex]
\hline

% Schritt 1: Goal abbauen (Allquantor und Implikation)
1. [Givens wie oben] \newline
3. Sei $\mathcal{A}$ ein beliebiges Mengensystem. \newline
4. $\mathcal{A} \subseteq \mathcal{T}_{\cap}$ (Annahme) &
\emph{Neues Goal:} $\bigcup \mathcal{A} \in \mathcal{T}_{\cap}$ \newline
\newline
(\emph{Umgewandelt via Definition von $\mathcal{T}_{\cap}$}) \newline
\emph{Goal*:} $\forall i \in I \, (\bigcup \mathcal{A} \in \mathcal{T}_i)$ \\[1ex]
\hline

% Schritt 2: Allquantor im Goal abbauen
1. [Givens wie oben]; 3. und 4. wie oben \newline
5. Sei $i_0 \in I$ ein beliebiges Indexelement. &
\emph{Neues Goal:} $\bigcup \mathcal{A} \in \mathcal{T}_{i_0}$ \\[1ex]
\hline

% Schritt 3: Das "Auge öffnen" bei Given 4
% Was bedeutet A \subseteq T_cap wirklich? Jedes Element aus A ist in JEDEM T_i.
1. [Givens wie oben]; 3. und 5. wie oben \newline
6. $\forall V \, (V \in \mathcal{A} \rightarrow V \in \mathcal{T}_{\cap})$ (Ausschreiben von 4) \newline
7. $\forall V \, \big( V \in \mathcal{A} \rightarrow \forall i \in I \, (V \in \mathcal{T}_i) \big)$ (Definition eingesetzt) &
\emph{Goal:} $\bigcup \mathcal{A} \in \mathcal{T}_{i_0}$ \\[1ex]
\hline

% Schritt 4: Given 7 für unser konkretes i_0 spezialisieren
1. [Givens wie oben]; 3. und 5. wie oben \newline
8. $\forall V \, (V \in \mathcal{A} \rightarrow V \in \mathcal{T}_{i_0})$ \newline
(\emph{Allquantor aus Given 7 für $i = i_0$ instanziiert}) \newline
\newline
\footnotesize{\textbf{Zwischenerkenntnis:}} \newline
\footnotesize{Given 8 bedeutet nichts anderes als: $\mathcal{A} \subseteq \mathcal{T}_{i_0}$!} &
\emph{Goal:} $\bigcup \mathcal{A} \in \mathcal{T}_{i_0}$ \\[1ex]
\hline

% Schritt 5: Axiom T3 der Topologie T_{i_0} zünden!
1. [Givens wie oben]; 3. und 5. wie oben \newline
9. $\mathcal{A} \subseteq \mathcal{T}_{i_0}$ (aus Schritt 4) \newline
10. $\mathcal{T}_{i_0}$ erfüllt T3, es gilt also: \newline
\phantom{XX} $\mathcal{A} \subseteq \mathcal{T}_{i_0} \rightarrow \bigcup \mathcal{A} \in \mathcal{T}_{i_0}$ \newline
11. $\bigcup \mathcal{A} \in \mathcal{T}_{i_0}$ (Modus Ponens aus 9 und 10) &
\emph{Goal:} $\bigcup \mathcal{A} \in \mathcal{T}_{i_0}$ \newline
\newline
\checkmark \quad (Ziel erreicht!) \\
\bottomrule
\end{tabular}
\end{center}

\end{document}
